\section{Problem Analysis}
\label{sec:problem-analysis}

Figure~\ref{fig:setting} shows the target setting, municipal rescue-service incident
response; the regional core also serves dispatch and reference data.

\subsection{Operational Requirements}
\label{sec:requirements}

R1 and R2 follow from mobile-IT and incident-reporting findings that responders need
locally available information and durable later
documentation~\cite{rsg-kommunalplan,pilemalm2014enabling,westman2021mobilisation}; R3
follows from decision-provenance requirements and rescue-service command
guidance~\cite{naja2010provenance,rsg-rad201}:
\begin{itemize}
  \item \textbf{\textit{R1 (Offline operation)}} Each tier must keep working from
  previously synchronized data without continuous connectivity.
  \item \textbf{\textit{R2 (Durable forward-only submission)}} Field submissions entered
  on tablets must survive vehicle restart and prolonged partition.
  \item \textbf{\textit{R3 (Attributable command authority)}} Each decision must be
  attributable to the one command authority in charge when it was made, and all
  decisions must form one order.
\end{itemize}

R3 needs authority and sequencing, but only for decisions. R1 and R2 need replication:
each node works from its own copy, and field submissions are stored locally and
forwarded later. Reference data, field submissions, and materialized views can therefore
use weaker, eventual contracts (\S~\ref{sec:state-taxonomy}).

Authority in crisis governance may be shared, delegated, contested, or legally bounded;
R3 abstracts only the accountability that decision provenance
requires~\cite{naja2010provenance}. We keep three things apart. \emph{Command authority}
belongs to a person, the incident commander, designated by the rescue-service
organization. The \emph{command role} is where the protocol exercises that authority:
one vehicle edge node, the command edge, identified by its node identity and current
epoch. \emph{Promotion} moves the role to another node. Only an authorized person, the
\emph{operator}---the commander, or command staff acting for the commander---can request
it, and the core then binds the next epoch to that node (\S~\ref{sec:protocol-m4}). The
role is never inferred from reachability.

\subsection{The Authority Trade-off under Partition}
\label{sec:tradeoff}

Consider a command edge A that loses contact with every other node. A can receive no
message, including one saying that command has moved. While A is cut off, a design can
provide at most two of the following:
\begin{enumerate}[label=(\alph*)]
  \item A keeps recording decisions;
  \item command moves to another vehicle, B, without A's participation;
  \item at most one vehicle records decisions at any moment.
\end{enumerate}
With (a) and (b), A and B both record. With (b) and (c), A must stop before B starts;
since A cannot be told, it must stop on a local timer, which bounds (a). With (a) and
(c), B must wait for A. This is the availability--consistency tension of partitioned
networks~\cite{davidson1985consistency}, applied to the authority role rather than to
the data. Table~\ref{tab:related-work} places each system family. Primary/backup systems
escape the trade-off with an out-of-band fence that powers the old primary
off~\cite{kleppmann2017ddia}; a vehicle out of radio range offers no such channel.
\systemname{} therefore has to choose, and it states its choice
(\S~\ref{sec:protocol-m4}): (a) for one lease window, (b) as soon as B reaches the core,
and (c) at the core but not at the edges.

\subsection{Positioning Relative to Consensus, CRDTs, Local-first, and Emergency Networking}
\label{sec:positioning}

\systemname{} recombines known mechanisms---single-leader replication, a transactional
outbox~\cite{richardson2018microservices}, idempotency keys, and fencing---around one
domain rule: the command edge is the only sequencer of decisions.

\textit{Consensus and replicated state machines.} These provide totally ordered logs
under their stated failure
models~\cite{lamport1998parttime,ongaro2014raft,balakrishnan2013tango,vanrenesse2004chain},
but elections select the writer by reachability, whereas R3 fixes it by operational
command (\S~\ref{sec:eval-raft}). A commander on the minority side cannot commit.

\textit{CRDTs.} Mergeable replicated
state~\cite{almeida2014delta,hanssen2025emergency,almeida2024crdtsurvey,mao2024reliablecrdt}
serves R1 and R2 well. As an authority mechanism, however, a multi-writer CRDT lets every
replica decide. When two responders set the same status to conflicting values, its merge
rule picks a winner, and no merge rule is defensible for such decisions. This rules
out merge as the source of authority, not CRDTs as replication. Once the command edge has
authorized and numbered its decisions, any eventually consistent mechanism can replicate
the log without conflicts: a grow-only set keyed by (epoch, \textit{event\_seq}) is a
CRDT in which each key has exactly one writer.

\textit{Offline-first and local-first systems.} These prioritize local
availability~\cite{kleppmann2019localfirst,richardson2018microservices} and suit
reference data, field submissions, and materialized views, but their goal is replica
\emph{autonomy}, which the command hierarchy rejects for decisions (R3).

\textit{Why not a centralized backend?} No server can order writes it cannot reach, and
partition from the incident is common. A
backend-authoritative design must either block decisions while the backhaul is down or
let field nodes write provisionally and merge later, which R3 forbids. \systemname{}
places the sequencer at the scene, so decisions continue through backhaul loss---for one
lease window, whose length sets how long (\S~\ref{sec:protocol-m4}).

\textit{Emergency and edge networking.} Closest in setting,
CoNICE~\cite{jahanian2022conice} reaches consensus among infrastructure-less peers using
One-Third-Rule agreement over epidemic replication; its no-privileged-peer assumption is
exactly the one R3 removes. Public-safety and edge/fog research places resilient
transport and storage near
responders~\cite{kumbhar2015legacy,wang2023overview,sagor2024distressnetng}, which serves
R1 but not the question of who may issue a decision.

\begin{table*}[t]
\centering
\caption{Where each system family sits in the authority trade-off of
\S~\ref{sec:tradeoff}, while the node that orders writes is cut off. At most two of
(a)--(c) can hold; only \systemname{} takes the writer from the command hierarchy.}
\label{tab:related-work}
\scriptsize
\setlength{\tabcolsep}{3pt}
\begin{tabularx}{\textwidth}{@{}>{\raggedright\arraybackslash}p{0.15\textwidth}>{\raggedright\arraybackslash}p{0.21\textwidth}>{\raggedright\arraybackslash}p{0.17\textwidth}>{\raggedright\arraybackslash}p{0.19\textwidth}>{\raggedright\arraybackslash}X@{}}
\toprule
\textbf{System} & \textbf{Writer chosen by} & \textbf{(a) Cut-off writer keeps recording} & \textbf{(b) Command moves without it} & \textbf{(c) One writer at a time} \\
\midrule
\textbf{\systemname{} (this paper)} & The command hierarchy: operator promotion, core-issued epoch & For one lease window & At once, if the new edge reaches the core & At the core; edges may overlap for one lease window \\
\addlinespace
Primary/backup with fencing~\cite{kleppmann2017ddia} & Configuration; manual failover & Until fenced & After an out-of-band fence & Yes \\
\addlinespace
Raft / shared log~\cite{ongaro2014raft,balakrishnan2013tango} & Election by the reachable quorum & Not on the minority side & Yes, by election & Yes \\
\addlinespace
CoNICE~\cite{jahanian2022conice} & Agreement among reachable peers & Only with peers that agree & Yes, by peer agreement & Yes \\
\addlinespace
RedBlue~\cite{li2012redblue} & Red: one global order; blue: any site & Blue operations only & No authority concept & For red operations \\
\addlinespace
Emergency CRDT-LWW~\cite{hanssen2025emergency} & Any replica & Yes & Yes & No: the merge picks the winner \\
\addlinespace
Local-first~\cite{kleppmann2019localfirst} & Any replica & Yes & Yes & No: replicas are autonomous \\
\bottomrule
\end{tabularx}
\end{table*}
